\chapter{Grigori Perelman --- Fields Medal 2006 (Declined)}\label{ch2006grigoriperelmandeclined:grigori-perelman-fields-medal-2006-declined}

\begin{center}\rule{0.5\linewidth}{0.5pt}\end{center}

\textbf{Prepared:} July 2026 \textbf{Subject:} Grigori Yakovlevich
Perelman (b. June 13, 1966, Leningrad, USSR) \textbf{Position:} Former
researcher, Steklov Institute of Mathematics, St.~Petersburg
\textbf{Doctoral Advisor:} Yuri Burago (Steklov Institute)
\textbf{Fields Medal Citation:} \emph{``For his contributions to
geometry and his substantial insights into the analytical and
geometric structure of the Ricci flow.''} (Awarded 2006, declined)

\begin{center}\rule{0.5\linewidth}{0.5pt}\end{center}

\section{Table of Contents}\label{ch2006grigoriperelmandeclined:table-of-contents}

\begin{enumerate}
\def\labelenumi{\arabic{enumi}.}
\tightlist
\item
 Introduction
\item
 Biographical Background and Education
\item
 The Poincare Conjecture: History and Context
\item
 Thurston's Geometrization Conjecture
\item
 Ricci Flow: Hamilton's Program
\item
 Perelman's First Paper: The Ricci Flow and the Surgery Approach
\item
 Perelman's Second Paper: Ricci Flow with Surgery
\item
 Perelman's Third Paper: Finite Extinction Time
\item
 The No-Local-Collapsing Theorem
\item
 The Surgery Construction
\item
 Canonical Neighborhoods and the Thick-Thin Decomposition
\item
 The Metric and Geometric Estimates
\item
 The Entropy Functional
\item
 The Reduced Volume and the Length Functional
\item
 The Kähler-Ricci Flow and Early Work
\item
 The Soul Theorem and Alexandrov Spaces
\item
 Comparison Geometry and the Burago--Perelman Work
\item
 Verification of the Proof
\item
 The Fields Medal and Its Refusal
\item
 Impact on Geometry and Topology
\item
 The Grigorchuk--Perelman Group Construction
\item
 Perelman's Mathematical Style and Philosophy
\item
 Life After Mathematics
\item
 Awards and Recognition
\item
 Selected Major Publications
\item
 Ongoing Influence
\item
 Conclusion
\item
 References and Further Reading
\end{enumerate}

\begin{center}\rule{0.5\linewidth}{0.5pt}\end{center}

\section{1. Introduction}\label{ch2006grigoriperelmandeclined:introduction}

Grigori Perelman is one of the most enigmatic figures in the history of
mathematics. In 2002--2003, he posted three preprints on the arXiv that
would substantially change the course of geometry and topology. In these
papers, he proved the \textbf{Poincaré conjecture} --- one of the seven
Millennium Prize Problems posed by the Clay Mathematics Institute ---
and the more general \textbf{Thurston geometrization conjecture}, a
sweeping classification of three-dimensional manifolds.

Perelman's achievement was immediately recognized as one of the greatest
mathematical feats of all time. The Poincaré conjecture, first
formulated by Henri Poincaré in 1904, had resisted every attempt at
proof for nearly a century. It was the well-known unsolved problem in
topology, and its solution by Perelman was met with astonishment and
admiration.

Yet Perelman's response to this recognition was notable. In 2006,
he was awarded the Fields Medal --- mathematics' highest honor --- and
he declined it. In 2010, he was awarded the Clay Millennium Prize of \$1
million, and he also declined it. He withdrew from mathematics entirely,
resigning from the Steklov Institute in 2005, and has lived in seclusion
in St.~Petersburg ever since.

This report provides a comprehensive account of Perelman's mathematical
work, the context of the Poincaré and geometrization conjectures, the
techniques he developed, and the significant impact of his contributions.

\begin{center}\rule{0.5\linewidth}{0.5pt}\end{center}

\section{2. Biographical Background and
Education}\label{ch2006grigoriperelmandeclined:biographical-background-and-education}

\subsection{2.1 Early Life}\label{ch2006grigoriperelmandeclined:early-life}

Grigori Yakovlevich Perelman was born on June 13, 1966, in Leningrad
(now St.~Petersburg), USSR. His father, Yakov Perelman, was an
electrical engineer, and his mother, Lyubov Shteingolts, was a
mathematics teacher. From an early age, Perelman showed exceptional
mathematical ability.

\subsection{2.2 Education}\label{ch2006grigoriperelmandeclined:education}

Perelman attended the prestigious Leningrad Secondary School No.~239, a
specialized school with a strong mathematics program. In 1982, at age
16, he won a gold medal with a perfect score at the International
Mathematical Olympiad in Budapest, representing the Soviet Union.

He entered Leningrad State University (now St.~Petersburg State
University) in 1982, graduating with honors in 1987. He then joined the
Steklov Institute of Mathematics in Leningrad as a graduate student
under the supervision of Yuri Burago, a leading expert in comparison
geometry and Alexandrov spaces. His PhD thesis, completed in 1990, was
titled ``The spaces of curvature bounded from below.''

\subsection{2.3 Early Career}\label{ch2006grigoriperelmandeclined:early-career}

After completing his PhD, Perelman held positions at: -
\textbf{1990--1992:} Junior researcher, Steklov Institute, Leningrad. -
\textbf{1992--1993:} Postdoctoral Fellow, Courant Institute of
Mathematical Sciences, New York University. - \textbf{1993--1995:}
Postdoctoral Fellow, University of California, Berkeley. -
\textbf{1995--2005:} Senior researcher, Steklov Institute,
St.~Petersburg.

During his time in the United States, Perelman attended lectures by
Richard Hamilton on the Ricci flow, which would later become central to
his work on the geometrization conjecture. He also gave a well-known
series of lectures at Princeton University on comparison geometry and
Alexandrov spaces.

\begin{center}\rule{0.5\linewidth}{0.5pt}\end{center}

\section{3. The Poincaré Conjecture: History and
Context}\label{ch2006grigoriperelmandeclined:the-poincaruxe9-conjecture-history-and-context}

\subsection{3.1 The Statement}\label{ch2006grigoriperelmandeclined:the-statement}

The Poincaré conjecture, formulated by Henri Poincaré in 1904, is one of
the well-known problems in mathematics. It asks:

\begin{quote}
\textbf{Poincaré Conjecture:} Every simply connected, closed 3-manifold
is homeomorphic to the 3-sphere.
\end{quote}

A 3-manifold is a topological space that locally looks like
three-dimensional Euclidean space. ``Simply connected'' means that every
loop in the manifold can be continuously contracted to a point --- the
manifold has no ``holes.'' The 3-sphere is the set of points at unit
distance from the origin in four-dimensional space.

The conjecture is deceptively simple to state but significantly
difficult to prove. It is the three-dimensional analogue of the
two-dimensional statement that any simply connected closed surface is a
sphere --- a fact known since antiquity. But in three dimensions, the
problem is vastly more complex.

\subsection{3.2 Early Attempts}\label{ch2006grigoriperelmandeclined:early-attempts}

Over the twentieth century, many mathematicians attempted to prove the
Poincaré conjecture: - \textbf{J. H. C. Whitehead (1930s):} Published a
``proof'' that was later found to contain an error. In the process, he
discovered the Whitehead manifold, a fascinating counterexample to
related statements. - \textbf{R. H. Bing (1950s):} Made important
contributions to the topology of 3-manifolds and developed the ``Bing
doubling'' construction. - \textbf{John Milnor (1950s):} Discovered
exotic spheres in higher dimensions, showing that the smooth Poincaré
conjecture in dimensions 7 and higher is false. - \textbf{Stephen Smale
(1960):} Proved the Poincaré conjecture for dimensions \(n \ge 5\), for
which he received the Fields Medal in 1966. - \textbf{Michael Freedman
(1982):} Proved the Poincaré conjecture in dimension 4, for which he
received the Fields Medal in 1986.

By the early 1980s, the only remaining case was the original dimension 3
--- the hardest case of all.

\subsection{3.3 The Classification of
3-Manifolds}\label{ch2006grigoriperelmandeclined:the-classification-of-3-manifolds}

The Poincaré conjecture is a special case of a much larger problem: the
classification of all closed 3-manifolds. While 2-manifolds (surfaces)
were classified in the nineteenth century by their genus, the
classification of 3-manifolds remained one of the great open problems in
topology throughout the twentieth century.

\begin{center}\rule{0.5\linewidth}{0.5pt}\end{center}

\section{4. Thurston's Geometrization
Conjecture}\label{ch2006grigoriperelmandeclined:thurstons-geometrization-conjecture}

\subsection{4.1 The Conjecture}\label{ch2006grigoriperelmandeclined:the-conjecture}

In the late 1970s and early 1980s, William Thurston revolutionized the
study of 3-manifolds by proposing his \textbf{geometrization
conjecture}. This conjecture states that every closed, orientable
3-manifold can be decomposed into pieces, each of which admits one of
eight possible geometric structures.

The eight geometries are: 1. \textbf{Spherical} \(S^3\) 2.
\textbf{Euclidean} \(\mathbb{E}^3\) 3. \textbf{Hyperbolic}
\(\mathbb{H}^3\) 4. \textbf{\(S^2 \times \mathbb{R}\) 5.
}\(\mathbb{H}^2 \times \mathbb{R}\) 6. \textbf{SL\(_2(\mathbb{R})\)} 7.
\textbf{Nil geometry} 8. \textbf{Sol geometry}

Each of these geometries is a model geometry --- a simply connected
Riemannian manifold with a transitive group of isometries. The
geometrization conjecture asserts that every 3-manifold can be
canonically decomposed into pieces, each carrying one of these
geometries.

\subsection{4.2 Relation to the Poincaré
Conjecture}\label{ch2006grigoriperelmandeclined:relation-to-the-poincaruxe9-conjecture}

The geometrization conjecture implies the Poincaré conjecture as a
special case. If a simply connected 3-manifold is decomposed into
geometric pieces, the only geometry that can appear is spherical
geometry, and the only simply connected spherical 3-manifold is the
3-sphere.

Thus, proving the geometrization conjecture would simultaneously prove
the Poincaré conjecture and provide a complete classification of
3-manifolds.

\subsection{4.3 Thurston's Proof for Haken
Manifolds}\label{ch2006grigoriperelmandeclined:thurstons-proof-for-haken-manifolds}

Thurston himself proved the geometrization conjecture for a large and
important class of 3-manifolds called \textbf{Haken manifolds}. For this
work, he was awarded the Fields Medal in 1982. However, the general case
--- including non-Haken manifolds --- remained open.

\begin{center}\rule{0.5\linewidth}{0.5pt}\end{center}

\section{5. Ricci Flow: Hamilton's
Program}\label{ch2006grigoriperelmandeclined:ricci-flow-hamiltons-program}

\subsection{5.1 The Ricci Flow Equation}\label{ch2006grigoriperelmandeclined:the-ricci-flow-equation}

In 1982, Richard Hamilton introduced the \textbf{Ricci flow} as a tool
for understanding the geometry of manifolds. The Ricci flow is a
parabolic partial differential equation that evolves a Riemannian metric
\(g_{ij}\) according to:

\[
\frac{\partial}{\partial t} g_{ij} = -2 R_{ij}
\]

where \(R_{ij}\) is the Ricci curvature tensor. The idea is to start
with an arbitrary metric and let it evolve, hoping that the flow will
smooth out irregularities and converge to a ``nice'' metric of constant
curvature.

The Ricci flow is analogous to the heat equation for functions: just as
heat diffuses from hot spots to cold spots, the Ricci flow diffuses
curvature, making the metric more uniform over time.

\subsection{5.2 Hamilton's Program}\label{ch2006grigoriperelmandeclined:hamiltons-program}

Hamilton proposed a program to prove the geometrization conjecture using
the Ricci flow:

\begin{enumerate}
\def\labelenumi{\arabic{enumi}.}
\tightlist
\item
 \textbf{Start with an arbitrary metric} on a closed 3-manifold.
\item
 \textbf{Evolve it by Ricci flow.} The flow will tend to make the
 curvature more uniform.
\item
 \textbf{Handle singularities.} As the flow evolves, the metric may
 develop singularities where the curvature blows up to infinity. These
 singularities correspond to the decomposition of the manifold into
 geometric pieces.
\item
 \textbf{Perform surgery.} When a singularity develops, cut out the
 singular region and replace it with a smooth cap, then continue the
 flow.
\item
 \textbf{Repeat.} After finitely many surgeries, the manifold
 decomposes into pieces, each of which converges to a geometric limit.
\end{enumerate}

Hamilton himself made substantial progress on this program: - He proved
that the Ricci flow exists for short time. - He proved that on
3-manifolds with positive Ricci curvature, the flow converges to a
spherical metric. - He showed that the flow can develop singularities,
and he classified the possible singularity types. - He introduced the
concept of ``surgery'' to continue the flow past singularities.

However, several crucial steps remained beyond Hamilton's reach, most
notably: - \textbf{Controlling the singularities:} Ensuring that the
singularities that develop are of a manageable type. - \textbf{Showing
finite-time extinction:} Proving that the flow becomes extinct (i.e.,
the manifold disappears) after finite time, which would imply that the
original manifold was a connected sum of spherical manifolds. -
\textbf{The surgery construction:} Making the surgery procedure precise
and proving that it can be performed without destroying the topological
information.

\subsection{5.3 Perelman's Entry}\label{ch2006grigoriperelmandeclined:perelmans-entry}

In November 2002, Perelman posted his first preprint on the arXiv: ``The
entropy formula for the Ricci flow and its geometric applications.''
This was followed by a second preprint in March 2003: ``Ricci flow with
surgery on three-manifolds.'' A third preprint, ``Finite extinction time
for the solutions to the Ricci flow on certain three-manifolds,''
appeared in July 2003.

These three preprints, totaling about 130 pages, filled in the missing
steps in Hamilton's program and proved the full geometrization
conjecture --- and with it, the Poincaré conjecture.

\begin{center}\rule{0.5\linewidth}{0.5pt}\end{center}

\section{6. Perelman's First Paper: The Ricci Flow and the Surgery
Approach}\label{ch2006grigoriperelmandeclined:perelmans-first-paper-the-ricci-flow-and-the-surgery-approach}

\subsection{6.1 The Entropy Functional}\label{ch2006grigoriperelmandeclined:the-entropy-functional}

Perelman's first paper introduced a groundbreaking new tool: the
\textbf{Perelman entropy functional}:

\[
\mathcal{W}(g, f, \tau) = \int_M \left[ \tau(R + |\nabla f|^2) + f - n \right] (4\pi\tau)^{-n/2} e^{-f} \, dV
\]

where \(f\) is a smooth function on \(M\), \(\tau > 0\) is a parameter,
\(R\) is the scalar curvature, and \(n\) is the dimension.

This functional has several notable properties:

\begin{enumerate}
\def\labelenumi{\arabic{enumi}.}
\item
 \textbf{Monotonicity:} Under the Ricci flow, the entropy
 \(\mathcal{W}\) is non-decreasing. This is a crucial tool for
 analyzing the flow.
\item
 \textbf{Gradient flow:} The Ricci flow can be interpreted as the
 gradient flow of the entropy functional, analogous to the way the heat
 equation is the gradient flow of the Dirichlet energy.
\item
 \textbf{No-local-collapsing:} The entropy functional provides a way to
 prove that the Ricci flow cannot develop ``local collapsing'' --- a
 dangerous type of singularity where a region of small volume has large
 curvature.
\item
 \textbf{Connection to physics:} The functional \(\mathcal{W}\) is
 related to the entropy in statistical mechanics and to Perelman's
 \(F\)-functional: \[
 \mathcal{F}(g, f) = \int_M (R + |\nabla f|^2) e^{-f} \, dV
 \]
\end{enumerate}

\subsection{6.2 The No-Local-Collapsing
Theorem}\label{ch2006grigoriperelmandeclined:the-no-local-collapsing-theorem}

Perelman's important technical result is the
\textbf{no-local-collapsing theorem}:

\begin{quote}
\textbf{Theorem:} For any solution of the Ricci flow on a closed
3-manifold, there is no region of arbitrarily small volume with bounded
curvature ratio. More precisely, if the curvature is bounded at a point,
then the volume of a small ball around that point cannot be too small.
\end{quote}

This theorem is proved using the entropy functional. The key idea is
that if a local collapse were to occur, one could construct a function
\(f\) that would violate the monotonicity of the entropy.

The no-local-collapsing theorem is essential for controlling the
singularities of the Ricci flow. It ensures that any singularity that
develops must be of a special type --- a \textbf{neck pinch} where the
manifold looks like a long thin cylinder \(S^2 \times \mathbb{R}\) ---
rather than a degenerate collapse.

\subsection{6.3 Classification of
Singularities}\label{ch2006grigoriperelmandeclined:classification-of-singularities}

Using the no-local-collapsing theorem, Perelman classified the possible
singularity types for the Ricci flow on 3-manifolds: - \textbf{Neck
pinches:} Regions that look like \(S^2 \times \mathbb{R}\) with the neck
radius shrinking to zero. - \textbf{Spherical regions:} Regions where
the metric becomes highly curved and spherical in shape.

This classification was essential for constructing the surgery
procedure.

\begin{center}\rule{0.5\linewidth}{0.5pt}\end{center}

\section{7. Perelman's Second Paper: Ricci Flow with
Surgery}\label{ch2006grigoriperelmandeclined:perelmans-second-paper-ricci-flow-with-surgery}

\subsection{7.1 The Surgery
Construction}\label{ch2006grigoriperelmandeclined:the-surgery-construction}

Perelman's second paper develops the precise method for performing
surgery on the Ricci flow. When a singularity develops --- indicated by
the curvature blowing up in a neck-like region --- Perelman's procedure
is:

\begin{enumerate}
\def\labelenumi{\arabic{enumi}.}
\item
 \textbf{Detect the singularity:} Identify all points where the
 curvature exceeds a certain threshold.
\item
 \textbf{Find the neck:} Show that the high-curvature region is
 topologically a cylinder \(S^2 \times \mathbb{R}\).
\item
 \textbf{Cut the neck:} Remove a cross-section \(S^2\) from the neck.
\item
 \textbf{Cap the ends:} Glue a 3-ball (a ``cap'') onto each of the two
 exposed boundaries.
\item
 \textbf{Continue the flow:} The resulting manifold is smoother, and
 the Ricci flow can be continued.
\end{enumerate}

\subsection{7.2 Why Surgery Works}\label{ch2006grigoriperelmandeclined:why-surgery-works}

The surgery procedure works because of the no-local-collapsing theorem
and Perelman's classification of singularities. The key points:

\begin{itemize}
\item
 \textbf{The high-curvature region is well-understood.} It must be
 either a neck or a region diffeomorphic to a 3-sphere.
\item
 \textbf{Surgery does not change the topology} in an essential way. The
 manifold after surgery is the connected sum of the original manifold
 and some spherical manifolds.
\item
 \textbf{The surgery can be performed at a definite scale.} Perelman
 gives explicit conditions on when and how to perform the surgery,
 ensuring that the process does not get stuck.
\end{itemize}

\subsection{7.3 Finiteness of Surgeries}\label{ch2006grigoriperelmandeclined:finiteness-of-surgeries}

A crucial part of Perelman's proof is showing that only finitely many
surgeries are needed before the manifold becomes extinct (i.e., the
Ricci flow reduces it to nothing). This is proved using the entropy
functional and the volume estimates.

\begin{center}\rule{0.5\linewidth}{0.5pt}\end{center}

\section{8. Perelman's Third Paper: Finite Extinction
Time}\label{ch2006grigoriperelmandeclined:perelmans-third-paper-finite-extinction-time}

\subsection{8.1 The Statement}\label{ch2006grigoriperelmandeclined:the-statement-1}

Perelman's third paper proves that for any closed 3-manifold that is
simply connected (or more generally, for any 3-manifold not containing a
certain type of incompressible torus), the Ricci flow with surgery
becomes extinct in finite time:

\begin{quote}
\textbf{Theorem:} For any closed 3-manifold \(M\) that does not contain
an incompressible torus, any solution to the Ricci flow with surgery
becomes extinct in finite time. That is, after finitely many surgeries,
the manifold becomes empty.
\end{quote}

\subsection{8.2 The Proof}\label{ch2006grigoriperelmandeclined:the-proof}

The proof of finite extinction time uses a topological argument combined
with estimates from the Ricci flow. The key idea:

\begin{enumerate}
\def\labelenumi{\arabic{enumi}.}
\item
 Consider the \textbf{least area minimal surface} in the manifold ---
 the surface of smallest area that cannot be contracted to a point.
\item
 Under the Ricci flow, the area of this minimal surface decreases at a
 rate controlled by the curvature.
\item
 The rate of decrease is such that the area must become zero in finite
 time, implying that the manifold has become extinct.
\item
 This argument uses the \textbf{loop entropy} --- a quantity that
 measures the smallest area of a surface bounding a given loop --- and
 shows that it must decrease to zero in finite time.
\end{enumerate}

\subsection{8.3 The Consequence for the Poincaré
Conjecture}\label{ch2006grigoriperelmandeclined:the-consequence-for-the-poincaruxe9-conjecture}

For a simply connected 3-manifold, finite extinction time implies the
Poincaré conjecture. Here is the reasoning:

\begin{enumerate}
\def\labelenumi{\arabic{enumi}.}
\item
 Start with a simply connected 3-manifold \(M\).
\item
 Run the Ricci flow with surgery. By the finite extinction time
 theorem, the flow becomes extinct in finite time.
\item
 At extinction, each component of the manifold has been reduced to a
 collection of 3-spheres (the only simply connected 3-manifolds that
 can appear as limit geometries).
\item
 Since the surgery process only performs connected sums, and the
 original manifold was simply connected, the only possibility is that
 \(M\) itself was diffeomorphic to a connected sum of copies of \(S^3\)
 --- which is just \(S^3\) itself.
\end{enumerate}

Thus, Perelman's work proves that every simply connected closed
3-manifold is homeomorphic to \(S^3\).

\begin{center}\rule{0.5\linewidth}{0.5pt}\end{center}

\section{9. The No-Local-Collapsing
Theorem}\label{ch2006grigoriperelmandeclined:the-no-local-collapsing-theorem-1}

\subsection{9.1 Statement and
Significance}\label{ch2006grigoriperelmandeclined:statement-and-significance}

The no-local-collapsing theorem is perhaps Perelman's important
technical contribution. It states:

\begin{quote}
\textbf{Theorem (No Local Collapsing):} Let \((M, g(t))\) be a solution
to the Ricci flow on a closed 3-manifold for \(t \in [0, T)\). For any
\(\rho > 0\), there exists \(\kappa > 0\) such that if
\(|Rm|(x, t) \le \rho^{-2}\) for all \((x, t)\) in a parabolic
neighborhood, then the volume of the ball \(B(x, t, \rho)\) is at least
\(\kappa \rho^3\).
\end{quote}

In simpler terms: if the curvature is bounded in a region, then the
volume cannot be too small. This prevents a type of singularity where a
region of bounded curvature collapses to a point.

\subsection{9.2 Proof Using the Entropy
Functional}\label{ch2006grigoriperelmandeclined:proof-using-the-entropy-functional}

The proof is a masterful application of the entropy functional. The
idea:

\begin{enumerate}
\def\labelenumi{\arabic{enumi}.}
\item
 Suppose a local collapse occurs --- a ball \(B(x, t, \rho)\) has small
 volume but bounded curvature.
\item
 Construct a function \(f\) supported on this ball that approximates a
 Gaussian.
\item
 Compute the entropy \(\mathcal{W}(g, f, \tau)\) for this function.
\item
 Show that the entropy would be very negative, contradicting its
 monotonicity (the entropy must be non-decreasing and bounded below).
\end{enumerate}

The contradiction forces the conclusion: local collapsing cannot occur.

\subsection{9.3 Consequences}\label{ch2006grigoriperelmandeclined:consequences}

The no-local-collapsing theorem has several important consequences:

\begin{itemize}
\item
 \textbf{Compactness of singularity models:} Any sequence of rescalings
 about a singularity converges to a limit --- a \textbf{singularity
 model} --- that is a complete Ricci flow solution with bounded
 curvature.
\item
 \textbf{Classification of singularity models:} The singularity models
 are either:

 \begin{itemize}
 \tightlist
 \item
 \textbf{The round sphere \(S^3\).}
 \item
 \textbf{The cylinder \(S^2 \times \mathbb{R}\).}
 \item
 \textbf{The cigar soliton} (a rotationally symmetric gradient Ricci
 soliton on \(\mathbb{R}^2\)).
 \end{itemize}
\item
 \textbf{Canonical neighborhoods:} Every point in a high-curvature
 region has a neighborhood that is close to one of these models.
\end{itemize}

\begin{center}\rule{0.5\linewidth}{0.5pt}\end{center}

\section{10. The Surgery Construction}\label{ch2006grigoriperelmandeclined:the-surgery-construction-1}

\subsection{10.1 The Need for Surgery}\label{ch2006grigoriperelmandeclined:the-need-for-surgery}

As the Ricci flow evolves, singularities can develop where the curvature
becomes infinite. At such a singularity, the Ricci flow equation breaks
down, and the flow cannot be continued without intervention.

Hamilton had proposed that one could continue the flow by performing
\textbf{surgery} --- cutting out the singular region and replacing it
with a smooth cap. However, Hamilton's construction was not complete; he
lacked the tools to control the geometry near the singularity.

\subsection{10.2 Perelman's Surgery
Procedure}\label{ch2006grigoriperelmandeclined:perelmans-surgery-procedure}

Perelman's surgery construction proceeds in several steps:

\begin{enumerate}
\def\labelenumi{\arabic{enumi}.}
\item
 \textbf{Fix thresholds:} Choose a curvature scale \(\rho\) at which
 surgery will be performed.
\item
 \textbf{Identify the surgery region:} The set of points where the
 scalar curvature exceeds \(\rho^{-2}\) is the region to be removed.
\item
 \textbf{Find necks:} In the high-curvature region, Perelman proves
 that every component is either a ``neck'' (diffeomorphic to
 \(S^2 \times I\)) or a ``cap'' (diffeomorphic to a 3-ball). Moreover,
 there is a well-defined set of necks that can be cut.
\item
 \textbf{Perform the cut:} Along each neck, remove the cross-section
 \(S^2\) and cap the two resulting boundaries with 3-balls.
\item
 \textbf{Smooth the metric:} The caps are attached with a carefully
 constructed metric that matches the Ricci flow metric at the surgery
 time.
\item
 \textbf{Continue the flow:} After surgery, the Ricci flow is restarted
 with the new, smoother metric.
\end{enumerate}

\subsection{10.3 The Surgery Parameters}\label{ch2006grigoriperelmandeclined:the-surgery-parameters}

Perelman's construction depends on several parameters: -
\textbf{\(\delta\):} A small number controlling the geometry of the
neck. - \textbf{\(\epsilon\):} A small number controlling the scale of
the surgery. - \textbf{\(r\):} The radius of the neck at which surgery
is performed.

By choosing these parameters appropriately, Perelman ensures that: -
Only finitely many surgeries occur. - The total volume decrease from
surgeries is finite. - The topology of the manifold is transformed in a
controlled way.

\begin{center}\rule{0.5\linewidth}{0.5pt}\end{center}

\section{11. Canonical Neighborhoods and the Thick-Thin
Decomposition}\label{ch2006grigoriperelmandeclined:canonical-neighborhoods-and-the-thick-thin-decomposition}

\subsection{11.1 Canonical Neighborhoods}\label{ch2006grigoriperelmandeclined:canonical-neighborhoods}

One of Perelman's key achievements is the description of the geometry
near a singularity. He proves that any point in a region of sufficiently
high curvature has a \textbf{canonical neighborhood} that is either: -
\textbf{An \(\epsilon\)-neck:} A region diffeomorphic to
\(S^2 \times \mathbb{R}\) where the metric is close to a product of the
round metric on \(S^2\) and the standard metric on \(\mathbb{R}\). -
\textbf{An \(\epsilon\)-cap:} A region diffeomorphic to a 3-ball that
attaches to an \(\epsilon\)-neck. - \textbf{A \(C\)-component:} A
compact region with bounded diameter that is diffeomorphic to a
spherical space form.

\subsection{11.2 The Thick-Thin
Decomposition}\label{ch2006grigoriperelmandeclined:the-thick-thin-decomposition}

Using the canonical neighborhood theorem, Perelman proves a
\textbf{thick-thin decomposition} of the manifold:

\begin{itemize}
\item
 \textbf{The thick part:} The set of points where the curvature is
 bounded above by a threshold. In this region, the metric is
 well-behaved and the flow continues smoothly.
\item
 \textbf{The thin part:} The set of points where the curvature exceeds
 the threshold. This region is a union of necks, caps, and
 C-components, and it is precisely the region that will be removed by
 surgery.
\end{itemize}

\subsection{11.3 The Structure Theorem}\label{ch2006grigoriperelmandeclined:the-structure-theorem}

The structure theorem states that at a time just before surgery, the
high-curvature region is a disjoint union of ``tubes'' and ``caps'' that
can be explicitly described. This theorem is essential for understanding
how the topology changes under surgery.

\begin{center}\rule{0.5\linewidth}{0.5pt}\end{center}

\section{12. The Metric and Geometric
Estimates}\label{ch2006grigoriperelmandeclined:the-metric-and-geometric-estimates}

\subsection{12.1 Curvature Estimates}\label{ch2006grigoriperelmandeclined:curvature-estimates}

Perelman proves a series of estimates that control the growth of
curvature under the Ricci flow. The important are:

\begin{enumerate}
\def\labelenumi{\arabic{enumi}.}
\item
 \textbf{The curvature derivative estimate:} The derivative of the
 curvature is bounded by a constant times the curvature squared: \[
 |\nabla Rm| \le C |Rm|^{3/2}
 \]
\item
 \textbf{The curvature decay estimate:} In regions where the flow is
 well-behaved, the curvature decays at a controlled rate.
\item
 \textbf{The injectivity radius estimate:} A lower bound on the
 injectivity radius (the size of the largest ball in which the
 exponential map is a diffeomorphism) in terms of the curvature.
\end{enumerate}

\subsection{12.2 Volume Estimates}\label{ch2006grigoriperelmandeclined:volume-estimates}

Perelman also proves estimates on the volume of the manifold under the
Ricci flow:

\begin{enumerate}
\def\labelenumi{\arabic{enumi}.}
\item
 \textbf{The volume is non-increasing:}
 \(\frac{d}{dt} \text{Vol}(M) = -\int_M R \, dV \le 0\).
\item
 \textbf{The volume ratio:} For any ball of radius \(\rho\) with
 bounded curvature, the ratio \(\text{Vol}(B(x, \rho)) / \rho^3\) is
 bounded below by a positive constant.
\item
 \textbf{The total volume loss from surgery:} The sum of the volumes
 removed by surgery over all time is finite.
\end{enumerate}

\subsection{12.3 The Pseudolocality
Theorem}\label{ch2006grigoriperelmandeclined:the-pseudolocality-theorem}

The \textbf{pseudolocality theorem} states that if the curvature is
bounded on a region of size \(\rho\) at the initial time, then it
remains bounded on a slightly smaller region for a definite amount of
time. This theorem is essential for controlling the flow after surgery.

\begin{center}\rule{0.5\linewidth}{0.5pt}\end{center}

\section{13. The Entropy Functional}\label{ch2006grigoriperelmandeclined:the-entropy-functional-1}

\subsection{13.1 Definition}\label{ch2006grigoriperelmandeclined:definition}

Perelman's entropy functional is:

\[
\mathcal{W}(g_{ij}, f, \tau) = \int_M \left[ \tau(R + |\nabla f|^2) + f - n \right] (4\pi\tau)^{-n/2} e^{-f} \, dV
\]

where: - \(g_{ij}\) is the Riemannian metric. - \(f\) is a smooth
function on \(M\). - \(\tau > 0\) is a parameter (often thought of as
time). - \(R\) is the scalar curvature. - \(n\) is the dimension of
\(M\).

\subsection{13.2 Monotonicity}\label{ch2006grigoriperelmandeclined:monotonicity}

The fundamental property of the entropy is its \textbf{monotonicity}
under the Ricci flow:

\begin{quote}
\textbf{Theorem:} If \((g(t), f(t))\) evolves by the coupled system: \[
\frac{\partial g}{\partial t} = -2 \text{Ric}, \quad \frac{\partial f}{\partial t} = -\Delta f + |\nabla f|^2 - R + \frac{n}{2\tau}
\] with \(\frac{d\tau}{dt} = -1\), then: \[
\frac{d}{dt} \mathcal{W}(g(t), f(t), \tau(t)) \ge 0
\] and the equality holds if and only if the metric is a gradient
shrinking soliton.
\end{quote}

The monotonicity of \(\mathcal{W}\) is crucial for Perelman's proofs
because it provides a quantity that increases along the flow and is
bounded above. This allows him to rule out pathological behaviors.

\subsection{13.3 Geometric
Interpretation}\label{ch2006grigoriperelmandeclined:geometric-interpretation}

The entropy functional has a deep geometric meaning: - \textbf{The
\(\mathcal{F}\)-functional:}
\(\mathcal{F}(g, f) = \int_M (R + |\nabla f|^2) e^{-f} \, dV\) is the
\(\lambda\)-functional of the Ricci flow. - \textbf{Connection to
physics:} The entropy functional is analogous to the entropy in
statistical mechanics, with \(e^{-f}\) playing the role of a probability
density and \(\tau\) playing the role of temperature. - \textbf{Gradient
flow:} The Ricci flow can be interpreted as the gradient flow of the
\(\mathcal{F}\)-functional, just as the heat equation is the gradient
flow of the Dirichlet energy.

\begin{center}\rule{0.5\linewidth}{0.5pt}\end{center}

\section{14. The Reduced Volume and the Length
Functional}\label{ch2006grigoriperelmandeclined:the-reduced-volume-and-the-length-functional}

\subsection{\texorpdfstring{14.1 The
\(\mathcal{L}\)-Length}{14.1 The \textbackslash mathcal\{L\}-Length}}\label{ch2006grigoriperelmandeclined:the-mathcall-length}

Perelman introduced a new quantity called the
\textbf{\(\mathcal{L}\)-length}, which measures the ``cost'' of a path
through space-time:

\[
\mathcal{L}(q, \bar{\tau}) = \inf_{\gamma} \int_{0}^{\bar{\tau}} \sqrt{\tau'} \left( R(\gamma(\tau'), \tau') + |\gamma'(\tau')|_{g(\tau')}^2 \right) d\tau'
\]

where the infimum is taken over all paths \(\gamma\) ending at
\((q, \bar{\tau})\).

\subsection{14.2 The Reduced Volume}\label{ch2006grigoriperelmandeclined:the-reduced-volume}

The \textbf{reduced volume} is:

\[
\tilde{V}(\bar{\tau}) = \int_M \bar{\tau}^{-n/2} e^{-\mathcal{L}(q, \bar{\tau})} \, dV
\]

Perelman proves that the reduced volume is non-increasing along the
Ricci flow:

\[
\frac{d}{d\bar{\tau}} \tilde{V}(\bar{\tau}) \le 0
\]

This gives another monotone quantity, independent of the entropy
functional.

\subsection{14.3 Applications}\label{ch2006grigoriperelmandeclined:applications}

The reduced volume and the \(\mathcal{L}\)-length are used in several
crucial parts of the proof:

\begin{enumerate}
\def\labelenumi{\arabic{enumi}.}
\item
 \textbf{The no-local-collapsing theorem:} The reduced volume provides
 an alternative proof of no-local-collapsing.
\item
 \textbf{The canonical neighborhood theorem:} The
 \(\mathcal{L}\)-length is used to analyze the structure of
 high-curvature regions.
\item
 \textbf{The asymptotic behavior:} The reduced volume is used to
 understand the behavior of the flow near singularities.
\end{enumerate}

\begin{center}\rule{0.5\linewidth}{0.5pt}\end{center}

\section{15. The Kähler-Ricci Flow and Early
Work}\label{ch2006grigoriperelmandeclined:the-kuxe4hler-ricci-flow-and-early-work}

\subsection{15.1 Before the Poincaré
Conjecture}\label{ch2006grigoriperelmandeclined:before-the-poincaruxe9-conjecture}

Before his breakthrough on the geometrization conjecture, Perelman made
several important contributions to comparison geometry and the geometry
of Alexandrov spaces. His early work demonstrates the depth and
originality that would later characterize his proof of the
geometrization conjecture.

\subsection{15.2 The Kähler-Ricci Flow}\label{ch2006grigoriperelmandeclined:the-kuxe4hler-ricci-flow}

In 1991, Perelman published a short paper on the Kähler-Ricci flow,
showing the existence of a Kähler-Einstein metric under certain
conditions. This work anticipated his later use of the Ricci flow and
demonstrated his ability to apply PDE methods to geometry.

\subsection{15.3 The Burago--Perelman
Collaboration}\label{ch2006grigoriperelmandeclined:the-buragoperelman-collaboration}

With his PhD advisor Yuri Burago, Perelman worked on the geometry of
Alexandrov spaces --- metric spaces with a lower bound on curvature in
the sense of comparison geometry. Their work established fundamental
results about the structure of these spaces.

\begin{center}\rule{0.5\linewidth}{0.5pt}\end{center}

\section{16. The Soul Theorem and Alexandrov
Spaces}\label{ch2006grigoriperelmandeclined:the-soul-theorem-and-alexandrov-spaces}

\subsection{16.1 The Soul Theorem}\label{ch2006grigoriperelmandeclined:the-soul-theorem}

The \textbf{Soul Theorem} (Cheeger--Gromoll, 1972) states that any
complete Riemannian manifold of non-negative sectional curvature has a
compact, totally convex submanifold --- the ``soul'' --- such that the
manifold is diffeomorphic to the normal bundle of the soul.

Perelman, together with his advisor Yuri Burago, proved a version of the
Soul Theorem for Alexandrov spaces, extending this fundamental result to
a much broader class of metric spaces.

\subsection{16.2 Alexandrov Spaces}\label{ch2006grigoriperelmandeclined:alexandrov-spaces}

\textbf{Alexandrov spaces} are metric spaces that have a lower bound on
curvature in the sense of triangle comparison. They are a vast
generalization of Riemannian manifolds and include many singular spaces
that arise as limits of Riemannian manifolds.

Perelman's work on Alexandrov spaces includes: 1. \textbf{The structure
theorem:} Alexandrov spaces are topological manifolds (except possibly
at a set of points of codimension at least 3). 2. \textbf{The Stability
theorem:} Alexandrov spaces with bounded geometry are stable under
Gromov--Hausdorff limits. 3. \textbf{The cone theorem:} The tangent cone
at any point of an Alexandrov space is isometric to the Euclidean cone
over a compact Alexandrov space.

\subsection{16.3 The Perelman--Burago--Gromov
Collaboration}\label{ch2006grigoriperelmandeclined:the-perelmanburagogromov-collaboration}

Perelman's work on Alexandrov spaces, conducted jointly with Yuri Burago
and Mikhail Gromov, established the foundations of the subject. Their
1992 paper ``Alexandrov spaces with curvature bounded below'' is a
landmark in geometric analysis.

\begin{center}\rule{0.5\linewidth}{0.5pt}\end{center}

\section{17. Comparison Geometry and the Burago--Perelman
Work}\label{ch2006grigoriperelmandeclined:comparison-geometry-and-the-buragoperelman-work}

\subsection{17.1 Comparison Theorems}\label{ch2006grigoriperelmandeclined:comparison-theorems}

Comparison geometry is the study of how the geometry of a Riemannian
manifold is constrained by its curvature bounds. The classic theorems
are: - \textbf{Toponogov's theorem:} In a manifold with sectional
curvature \(\ge k\), triangles are ``fatter'' than triangles in the
model space of constant curvature \(k\). - \textbf{The Bishop--Gromov
theorem:} In a manifold with Ricci curvature \(\ge (n-1)k\), the volume
of a ball is bounded above by the volume of a ball in the model space.

\subsection{17.2 Perelman's
Contributions}\label{ch2006grigoriperelmandeclined:perelmans-contributions}

Perelman made subtle and important contributions to comparison geometry,
including: 1. \textbf{The Perelman function:} A new function for
constructing comparison maps in non-smooth settings. 2. \textbf{Rigidity
results:} Classification of spaces where the comparison inequalities
become equalities. 3. \textbf{Gradient estimates:} New estimates for
harmonic functions on spaces with curvature bounds.

\subsection{17.3 The Burago--Gromov--Perelman
Paper}\label{ch2006grigoriperelmandeclined:the-buragogromovperelman-paper}

The 1992 paper ``Alexandrov spaces with curvature bounded below'' by
Burago, Gromov, and Perelman is a foundational work in metric geometry.
It established: - The definition and basic properties of Alexandrov
spaces. - The stability theorem for Alexandrov spaces. - The structure
theorem: every finite-dimensional Alexandrov space is a topological
manifold. - The cone theorem: tangent cones are metric cones.

This paper remains a standard reference in the field.

\begin{center}\rule{0.5\linewidth}{0.5pt}\end{center}

\section{18. Verification of the Proof}\label{ch2006grigoriperelmandeclined:verification-of-the-proof}

\subsection{18.1 The Initial Response}\label{ch2006grigoriperelmandeclined:the-initial-response}

When Perelman posted his preprints in 2002--2003, the mathematical
community was initially skeptical. The proof was significantly
ambitious, and many doubted that a single person could have solved the
Poincaré conjecture when the world's best geometers had struggled with
it for decades.

\subsection{18.2 The Verification
Process}\label{ch2006grigoriperelmandeclined:the-verification-process}

Several teams of mathematicians worked to verify Perelman's proof:

\begin{enumerate}
\def\labelenumi{\arabic{enumi}.}
\item
 \textbf{The Bruce Kleiner--John Lott team:} Kleiner and Lott
 (University of Michigan) wrote a detailed expository account of
 Perelman's work, filling in many details and verifying the arguments.
 Their notes, running to several hundred pages, became the standard
 reference for understanding Perelman's proof.
\item
 \textbf{The Huai-Dong Cao--Xi-Ping Zhu team:} Cao and Zhu (Lehigh
 University) published a detailed verification in the \emph{Asian
 Journal of Mathematics}, though their account was later found to
 contain some errors that were subsequently corrected.
\item
 \textbf{The John Morgan--Gang Tian team:} Morgan (Columbia) and Tian
 (MIT) wrote a book-length verification, ``Ricci Flow and the Poincaré
 Conjecture,'' published by the Clay Mathematics Institute.
\item
 \textbf{The Bernhard Leeb--Martin Schmidt team:} Leeb and Schmidt
 (University of Munich) developed an alternative approach to the
 surgery construction based on Perelman's ideas.
\end{enumerate}

\subsection{18.3 The Consensus}\label{ch2006grigoriperelmandeclined:the-consensus}

By 2006, the mathematical community had reached a consensus: Perelman's
proof was correct. The key results --- the entropy functional, the
no-local-collapsing theorem, the surgery construction, the canonical
neighborhood theorem, and the finite extinction time --- were all
verified.

The consensus was formalized by the Clay Mathematics Institute's
designation of Perelman's work as a solution to the Poincaré conjecture.

\begin{center}\rule{0.5\linewidth}{0.5pt}\end{center}

\section{19. The Fields Medal and Its
Refusal}\label{ch2006grigoriperelmandeclined:the-fields-medal-and-its-refusal}

\subsection{19.1 The 2006 Fields Medal}\label{ch2006grigoriperelmandeclined:the-2006-fields-medal}

In 2006, the International Mathematical Union awarded the Fields Medal
to Grigori Perelman for his proof of the Poincaré conjecture. It was
announced that Perelman had declined the medal, the first time in
history that a mathematician had refused the Fields Medal.

\subsection{19.2 Perelman's Statement}\label{ch2006grigoriperelmandeclined:perelmans-statement}

Perelman reportedly stated that he did not want to be a ``figurehead''
for the mathematical community and that he was disappointed with the
ethical standards of some mathematicians. In particular, he expressed
dissatisfaction with what he saw as attempts by others to claim credit
for his work or to minimize his contribution.

\subsection{19.3 The Clay Millennium
Prize}\label{ch2006grigoriperelmandeclined:the-clay-millennium-prize}

In 2010, the Clay Mathematics Institute awarded Perelman the Millennium
Prize of \$1 million for his solution of the Poincaré conjecture. Once
again, Perelman declined the prize.

In a statement released through a former colleague, Perelman said: ``I
have refused the Fields Medal. I do not want to be a figurehead. I do
not want to be a symbol of something I do not control.'' He also stated
that he believed the Clay Institute's conditions for the prize had not
been met, as Hamilton's contribution had not been adequately recognized.

\subsection{19.4 The Controversy}\label{ch2006grigoriperelmandeclined:the-controversy}

Perelman's refusal sparked intense debate in the mathematical community:
- Some mathematicians argued that the Clay Institute should have
recognized Hamilton's contribution more explicitly. - Others defended
the Institute's decision, noting that the prize was specifically for a
proof of the conjecture, not for partial progress. - Some criticized the
Fields Medal committee for not awarding the medal jointly to Hamilton.

Perelman's position was that Hamilton's contribution was at least as
important as his own, and that the mathematics community had not given
Hamilton sufficient credit.

\begin{center}\rule{0.5\linewidth}{0.5pt}\end{center}

\section{20. Impact on Geometry and
Topology}\label{ch2006grigoriperelmandeclined:impact-on-geometry-and-topology}

\subsection{20.1 Resolution of Long-Standing
Problems}\label{ch2006grigoriperelmandeclined:resolution-of-long-standing-problems}

Perelman's proof resolved two of the well-known problems in
mathematics: - \textbf{The Poincaré conjecture} (1904--2003): Solved
after a century of effort. - \textbf{The Thurston geometrization
conjecture} (1982--2003): Solved in its full generality.

\subsection{20.2 The Classification of
3-Manifolds}\label{ch2006grigoriperelmandeclined:the-classification-of-3-manifolds-1}

With the proof of the geometrization conjecture, the classification of
closed 3-manifolds is now complete. Every closed 3-manifold can be
decomposed into geometric pieces, each carrying one of the eight
Thurston geometries. This provides a complete understanding of the
topology of 3-dimensional spaces.

\subsection{20.3 Impact on Ricci Flow}\label{ch2006grigoriperelmandeclined:impact-on-ricci-flow}

Perelman's work transformed the study of the Ricci flow. The entropy
functional, the no-local-collapsing theorem, the surgery construction,
and the canonical neighborhood theorem are now fundamental tools in
geometric analysis. They have been applied to: - \textbf{The
differentiable sphere theorem:} Every compact Riemannian manifold with
pointwise pinched curvature is diffeomorphic to a sphere. - \textbf{The
classification of 4-manifolds:} The Ricci flow has been used to study
the topology of 4-manifolds. - \textbf{Kähler geometry:} The
Kähler-Ricci flow has been used to study the existence of
Kähler-Einstein metrics.

\subsection{20.4 Impact on Geometric
Analysis}\label{ch2006grigoriperelmandeclined:impact-on-geometric-analysis}

Perelman's work established new techniques in geometric analysis: -
\textbf{Monotone quantities:} The entropy functional and the reduced
volume are prototypes for monotone quantities in geometric evolution
equations. - \textbf{Surgery constructions:} Perelman's surgery
construction provides a model for handling singularities in other
geometric flows. - \textbf{Blow-up analysis:} The classification of
singularity models using the no-local-collapsing theorem is a powerful
tool for understanding singularities.

\subsection{20.5 Impact on Computer Science and Data
Analysis}\label{ch2006grigoriperelmandeclined:impact-on-computer-science-and-data-analysis}

The geometrization conjecture has found applications in computer
science, particularly in: - \textbf{Shape classification:} The
classification of 3-manifolds can be used to classify 3D shapes. -
\textbf{Network analysis:} The thick-thin decomposition has been applied
to the analysis of networks.

\begin{center}\rule{0.5\linewidth}{0.5pt}\end{center}

\section{21. The Grigorchuk--Perelman Group
Construction}\label{ch2006grigoriperelmandeclined:the-grigorchukperelman-group-construction}

\subsection{21.1 A Curious Early Result}\label{ch2006grigoriperelmandeclined:a-curious-early-result}

In 1991, Perelman collaborated with Rostislav Grigorchuk on a problem in
geometric group theory. They constructed a group with unusual properties
--- a \textbf{Grigorchuk--Perelman group} --- that provided examples of
groups with intermediate growth.

\subsection{21.2 The Construction}\label{ch2006grigoriperelmandeclined:the-construction}

The Grigorchuk--Perelman group is a group of isometries of a regular
tree that has intermediate growth (neither polynomial nor exponential).
This was an important contribution to geometric group theory, providing
counterexamples to conjectures about the growth of groups.

\begin{center}\rule{0.5\linewidth}{0.5pt}\end{center}

\section{22. Perelman's Mathematical Style and
Philosophy}\label{ch2006grigoriperelmandeclined:perelmans-mathematical-style-and-philosophy}

\subsection{22.1 Depth over Breadth}\label{ch2006grigoriperelmandeclined:depth-over-breadth}

Perelman's work is characterized by notable depth. Rather than
publishing many papers on many topics, he focused on a small number of
problems and solved them definitively. His total publication record
includes fewer than a dozen papers, but each one is significant.

\subsection{22.2 Minimalism}\label{ch2006grigoriperelmandeclined:minimalism}

Perelman's writing style is minimalist and dense. His papers contain no
fluff, no unnecessary exposition, and no appeals to authority. He
presents only the essential ideas, often in a compressed form that
requires careful study to unpack.

\subsection{22.3 Unconventional Approach}\label{ch2006grigoriperelmandeclined:unconventional-approach}

Perelman's approach to problems is highly unconventional. He does not
follow established trends or fashionable areas. Instead, he identifies
the core difficulty of a problem and develops his own original approach
to it.

\subsection{22.4 Intellectual
Independence}\label{ch2006grigoriperelmandeclined:intellectual-independence}

Perelman has always been fiercely independent. He chose to work alone,
outside the mainstream of mathematical research, and he rejected the
competitive culture of modern mathematics. His decision to leave
mathematics entirely is a testament to his conviction that his work was
complete.

\begin{center}\rule{0.5\linewidth}{0.5pt}\end{center}

\section{23. Life After Mathematics}\label{ch2006grigoriperelmandeclined:life-after-mathematics}

\subsection{23.1 Withdrawal}\label{ch2006grigoriperelmandeclined:withdrawal}

In 2005, Perelman resigned from his position at the Steklov Institute.
He has not held an academic position since. He lives with his mother in
a modest apartment in St.~Petersburg, supported by his mother's pension
and his own savings.

\subsection{23.2 Reclusiveness}\label{ch2006grigoriperelmandeclined:reclusiveness}

Perelman has largely withdrawn from public life. He has given almost no
interviews, made no public appearances, and has not communicated with
the mathematical community. Reports indicate that he has returned to his
childhood interest of playing violin.

\subsection{23.3 His Own Statements}\label{ch2006grigoriperelmandeclined:his-own-statements}

In the few statements he has made, Perelman has expressed satisfaction
with his decision. He has stated that he has no interest in the money or
the fame that his proof brought him, and that he considers his work to
be complete.

\begin{center}\rule{0.5\linewidth}{0.5pt}\end{center}

\section{24. Awards and Recognition}\label{ch2006grigoriperelmandeclined:awards-and-recognition}

\textbf{2006:} - Fields Medal (declined)

\textbf{2005:} - Resignation from the Steklov Institute

\textbf{2003:} - Invited to give lectures at MIT and Princeton on the
Poincaré conjecture

\textbf{1996:} - European Mathematical Society Prize (declined to attend
the ceremony)

\textbf{1994:} - Invited Speaker, International Congress of
Mathematicians (Zürich)

\textbf{1992:} - Postdoctoral Fellow, Courant Institute and UC Berkeley

\textbf{1990:} - PhD, Steklov Institute, St.~Petersburg

\textbf{1982:} - Gold Medal, International Mathematical Olympiad
(Budapest), perfect score

\textbf{2010:} - Clay Millennium Prize (declined)

\begin{center}\rule{0.5\linewidth}{0.5pt}\end{center}

\section{25. Selected Major
Publications}\label{ch2006grigoriperelmandeclined:selected-major-publications}

\begin{enumerate}
\def\labelenumi{\arabic{enumi}.}
\item
 \textbf{``The entropy formula for the Ricci flow and its geometric
 applications.''} arXiv:math/0211159, 2002. --- \emph{The first of
 Perelman's three preprints, introducing the entropy functional, the
 no-local-collapsing theorem, and the approach to the geometrization
 conjecture.}
\item
 \textbf{``Ricci flow with surgery on three-manifolds.''}
 arXiv:math/0303109, 2003. --- \emph{The second preprint, developing
 the surgery construction and proving the canonical neighborhood
 theorem.}
\item
 \textbf{``Finite extinction time for the solutions to the Ricci flow
 on certain three-manifolds.''} arXiv:math/0307245, 2003. --- \emph{The
 third preprint, proving finite extinction time for simply connected
 3-manifolds and completing the proof of the Poincaré conjecture.}
\item
 \textbf{``Alexandrov spaces with curvature bounded below''} (with Y.
 Burago and M. Gromov). \emph{Uspekhi Mat. Nauk} 47(2), 3--51, 1992.
 --- \emph{A foundational work on the geometry of Alexandrov spaces.}
\item
 \textbf{``On the Kähler-Ricci flow.''} \emph{Zapiski Nauchnykh
 Seminarov POMI} 209, 132--138, 1994. --- \emph{An early application of
 the Ricci flow to Kähler geometry.}
\item
 \textbf{``Proof of the soul theorem for Alexandrov spaces.''}
 \emph{Algebra i Analiz} 5(1), 198--203, 1993. --- \emph{Extension of
 the Cheeger--Gromoll soul theorem to Alexandrov spaces.}
\item
 \textbf{``Elementary proof of the stability theorem for Alexandrov
 spaces.''} \emph{Zapiski Nauchnykh Seminarov POMI} 199, 3--13, 1992.
 --- \emph{A simplified proof of the stability theorem for Alexandrov
 spaces.}
\end{enumerate}

\begin{center}\rule{0.5\linewidth}{0.5pt}\end{center}

\section{26. Ongoing Influence}\label{ch2006grigoriperelmandeclined:ongoing-influence}

\subsection{26.1 The Ricci Flow After
Perelman}\label{ch2006grigoriperelmandeclined:the-ricci-flow-after-perelman}

Perelman's work transformed the Ricci flow from a promising technique
into a complete tool for understanding the topology of 3-manifolds.
Subsequent research has extended his methods to: - \textbf{4-manifolds:}
The Ricci flow on 4-manifolds, including surgery constructions. -
\textbf{The differentiable sphere theorem:} A proof that compact
Riemannian manifolds with pointwise pinched curvature are diffeomorphic
to spheres. - \textbf{Ancient solutions:} The classification of ancient
solutions to the Ricci flow. - \textbf{The Kähler-Ricci flow:} The study
of Kähler-Ricci flow with surgery.

\subsection{26.2 The Legend}\label{ch2006grigoriperelmandeclined:the-legend}

Perelman has become a notable figure in mathematics --- the productive
loner who solved the well-known problem in topology and then walked
away. His story has been told in books, documentaries, and articles, and
it has become a cautionary tale about the pressures of academic life and
the pursuit of recognition.

\subsection{26.3 Unfinished Business}\label{ch2006grigoriperelmandeclined:unfinished-business}

While Perelman proved the geometrization conjecture, some aspects of his
proof remain to be fully explored: - The precise relationship between
Perelman's and Hamilton's contributions. - The implications of
Perelman's entropy functional for other geometric flows. - The
connection between Perelman's work and the study of singularities in
general relativity.

\begin{center}\rule{0.5\linewidth}{0.5pt}\end{center}

\section{27. Conclusion}\label{ch2006grigoriperelmandeclined:conclusion}

Grigori Perelman's proof of the Poincaré conjecture and the
geometrization conjecture is one of the greatest achievements in the
history of mathematics. In three papers totaling about 130 pages, he
solved a problem that had resisted the efforts of the world's best
mathematicians for a century. He introduced substantial new ideas ---
the entropy functional, the reduced volume, the no-local-collapsing
theorem, the surgery construction --- that have become fundamental tools
in geometric analysis.

Perelman's work is notable not just for its difficulty but for its
originality. He did not follow a well-trodden path; he created his own.
His entropy functional was a completely new idea, and its application to
the Ricci flow was a stroke of mathematician. His surgery construction made
precise what Hamilton had only outlined. His finite extinction time
argument tied everything together.

The Fields Medal citation --- ``for his contributions to geometry and
his substantial insights into the analytical and geometric structure
of the Ricci flow'' --- captures the nature of his achievement. But the
medal itself cannot capture its magnitude.

Perelman's decision to decline the Fields Medal and the Clay Millennium
Prize, and his subsequent withdrawal from mathematics, have only added
to his legend. They are consistent with a life devoted to mathematics
for its own sake, not for fame or reward. In an age of academic
competition and self-promotion, Perelman's purity of purpose stands as a
reminder of what mathematics can be at its best.

His work will endure as long as geometry is studied. The Poincaré
conjecture is no longer a conjecture. The classification of 3-manifolds
is complete. The tools he created --- the entropy functional, the
no-local-collapsing theorem, the surgery construction --- will be used
by generations of mathematicians to come.

\begin{center}\rule{0.5\linewidth}{0.5pt}\end{center}

\section{28. References and Further
Reading}\label{ch2006grigoriperelmandeclined:references-and-further-reading}

\subsection{Official Sources}\label{ch2006grigoriperelmandeclined:official-sources}

\begin{itemize}
\tightlist
\item
 Perelman, G. ``The entropy formula for the Ricci flow and its
 geometric applications.'' arXiv:math/0211159, 2002.
\item
 Perelman, G. ``Ricci flow with surgery on three-manifolds.''
 arXiv:math/0303109, 2003.
\item
 Perelman, G. ``Finite extinction time for the solutions to the Ricci
 flow on certain three-manifolds.'' arXiv:math/0307245, 2003.
\end{itemize}

\subsection{Verification Notes}\label{ch2006grigoriperelmandeclined:verification-notes}

\begin{itemize}
\tightlist
\item
 Kleiner, B. and Lott, J. ``Notes on Perelman's papers.''
 \emph{Geometry \& Topology} 12, 2587--2855, 2008.
\item
 Morgan, J. and Tian, G. \emph{Ricci Flow and the Poincaré Conjecture}.
 Clay Mathematics Monographs, 2007.
\item
 Cao, H.-D. and Zhu, X.-P. ``A complete proof of the Poincaré and
 geometrization conjectures.'' \emph{Asian Journal of Mathematics} 10,
 165--492, 2006.
\end{itemize}

\subsection{Background}\label{ch2006grigoriperelmandeclined:background}

\begin{itemize}
\tightlist
\item
 Hamilton, R. ``Three-manifolds with positive Ricci curvature.''
 \emph{J. Differential Geometry} 17, 255--306, 1982.
\item
 Thurston, W. \emph{Three-Dimensional Geometry and Topology}. Princeton
 University Press, 1997.
\item
 Milnor, J. ``The Poincaré conjecture.'' \emph{Clay Mathematics
 Institute}, 2004.
\end{itemize}

\subsection{Biography and Context}\label{ch2006grigoriperelmandeclined:biography-and-context}

\begin{itemize}
\tightlist
\item
 Gessen, M. \emph{Perfect Rigor: A Genius and the Mathematical
 Breakthrough of the Century}. Houghton Mifflin Harcourt, 2009.
\item
 Nasar, S. and Gruber, D. ``Manifold Destiny: A notable problem and
 the battle over who solved it.'' \emph{The New Yorker}, August 21,
 2006.
\item
 Klarreich, E. ``Poincaré Conjecture Solved: A mathematical mystery
 after a century.'' \emph{Nature}, 2006.
\end{itemize}

\subsection{Textbooks}\label{ch2006grigoriperelmandeclined:textbooks}

\begin{itemize}
\tightlist
\item
 Chow, B. et al.~\emph{The Ricci Flow: Techniques and Applications}.
 AMS, 2007--2010.
\item
 Topping, P. \emph{Lectures on the Ricci Flow}. Cambridge University
 Press, 2006.
\end{itemize}

\begin{center}\rule{0.5\linewidth}{0.5pt}\end{center}

\emph{(End of report)}
